% 390_Puzzleteile_Etabliert_En_ch.tex
% Johann Pascher, 9 Oct 2026
% Loose Puzzle Pieces: Kaluza-Klein and Other Established Building Blocks of the Time-Mass Duality

\providecommand{\BM}{\textbf{[B]}}
\providecommand{\KM}{\textbf{[K]}}
\providecommand{\SM}{\textbf{[S]}}
\providecommand{\XM}{\textbf{[X]}}
\providecommand{\QM}{\textbf{[Q]}}
\providecommand{\EM}{\textbf{[E]}}
\providecommand{\SetzM}{\textbf{[SETZUNG]}}

\chapter*{Loose Puzzle Pieces: Kaluza-Klein and Other Established Building Blocks of the Time-Mass Duality}
\addcontentsline{toc}{chapter}{Loose Puzzle Pieces: Kaluza-Klein and Other Established Building Blocks}

\begin{abstract}
Anyone who sees $\tilde T\cdot m=1$ for the first time often recognises something familiar in it,
most often Kaluza-Klein. That is not an objection but a finding: almost every building block of
FFGFT has been in the literature for decades, scattered, however, across relativity, quantum
mechanics, group theory, string theory, number theory and phenomenology. This document places
eleven such building blocks: Kaluza-Klein, de Broglie's internal clock and the Compton time, the
Zitterbewegung, Pauli's objection to a time operator, the relational approaches to time, Wigner's
mass as a Casimir invariant, the Casimir value of $SU(3)$, Heaviside-Lorentz units, orbifold
compactification, finite fields with their Frobenius automorphism, and the Koide formula. For each
building block it records what it already contained, which piece was missing, and where in the
corpus it fits in. The order matters: these building blocks were not the foundation of FFGFT. The
theory arose independently of them from the relation $\tilde T\cdot m=1$; the building blocks were
found in the literature only afterwards and fit into a picture that was drawn without them. The
result is a map of the places at which the theory could already have been found.
\end{abstract}

% ============================================================
\section*{Status markers}
% ============================================================
\begin{center}
\begin{tabular}{@{}lp{11cm}@{}}
\textbf{[SETZUNG]} & Motivated input, not a derivation. \\[2pt]
\textbf{[B]} & Algebraically proved (verification script, exact arithmetic). \\[2pt]
\textbf{[K]} & Numerically confirmed. \\[2pt]
\textbf{[S]} & Structurally plausible, identification open. \\[2pt]
\textbf{[E]} & Established mathematics or physics, adopted. \\[2pt]
\textbf{[Q]} & Taken from the literature, with source. \\[2pt]
\textbf{[X]} & Checked and negative. \\
\end{tabular}
\end{center}
The markers in this document give the status \emph{in the corpus}; the document itself derives
nothing new but places things in context. Verification script: \texttt{pruef\_390\_puzzleteile.py}.

% ============================================================
\section{The question}
% ============================================================

In the introductory conversation on the website, the sceptical physicist asks right at the start
whether this does not sound like Kaluza-Klein. The answer is yes, and the question can be
generalised: if the duality of mass and time was already known, why was it not thought through to
the end?

For Einstein, the corpus has already answered this question. Doc.~312, chapter \enquote{Why
Einstein did not see $\tilde T\cdot m=1$}, shows that the ingredients were on his desk and that the
missing step was not a formula but the change of status of a known formula into a principle. The
present document does not repeat that account. It widens the view: not just one ingredient but
almost all of them were available, each in a different field, each treated as a tool or a
curiosity. The following sections go through them in three groups: the rolled-up direction, mass
as structure, and the geometry with its algebra.

\textbf{On the order.} None of the following building blocks was a starting point of FFGFT. The
theory was developed from $\tilde T\cdot m=1$ and the geometry of the carrier; the literature
references were looked for and found only afterwards. Where it says below that a building block fits
in at some place in the corpus, this is to be read as agreement in hindsight, not as origin. This is
precisely what gives the map its value: pieces that arose independently of one another and
independently of FFGFT fit into the same picture.

% ============================================================
\section{Time and mass: the rolled-up direction}
% ============================================================

\subsection{Kaluza-Klein (1921, 1926)}

\textbf{What it contained.} Kaluza extended general relativity by a fifth dimension, and Klein
rolled it up into a circle of radius $R$ \QM~\cite{Kaluza1921,Klein1926}. On a circle only integer
windings are allowed; from this a mode tower with $m_n\propto n/R$ arises, and momentum along the
circle direction appears in four dimensions as mass. Klein thereby also obtained the quantisation
of charge. So by 1926 the core statement was in place that a compact direction generates mass as
the inverse of its period.

\textbf{What was missing.} The rolled-up direction was spatial and was added to spacetime. The
radius $R$ remained a free parameter, and the predicted tower of heavy excitations was never found.
As Doc.~311 puts it, one had gained a scale and bought a parameter along with it.

\textbf{Where it fits in the corpus.} In the corpus, $\tilde T\cdot m=1$ is explicitly the
Kaluza-Klein relation of the fourth torus direction, the mass dimension $S^1_m$ (Doc.~330,
Doc.~270). The direction is timelike: what holds there for a rolled-up spatial direction holds here
for rolled-up time, with mass in the role of the mode index (Doc.~307; likewise Doc.~231). The
radius is not a free modulus, because $\tilde T\cdot m=1$ binds the direction; Doc.~311 summarises
this as \enquote{scale without modulus}, and Doc.~318 sets the three routes side by side. There is
no additional tower to wait for: the discrete spectrum of the compact direction already is the
observed particle spectrum (Doc.~307). Locally, Lorentz invariance holds without restriction; the
compactness shows only globally (Doc.~312).

\subsection{De Broglie's internal clock and the Compton time (1924, 2013)}

\textbf{What it contained.} De Broglie assigned to every particle an internal clock with frequency
$\nu=mc^2/h$ \QM~\cite{deBroglie1925}. In natural units this is $T\cdot m=1$ in raw form. The
Compton time $\hbar/(mc^2)$ has been a standard notion ever since, and in 2013 Lan, Müller and
co-workers built a clock whose rate follows directly from the mass of an atom \QM~\cite{Lan2013}.

\textbf{What was missing.} The clock became the wavelength; the clock itself dropped out as a
derived scale. The history of this step is told in Doc.~312 and is not repeated here.

\textbf{Where it fits in the corpus.} The de Broglie relation appears on the mass direction as
$\lambda_4\cdot m=2\pi$, exact mode by mode (Doc.~314); the clock character of every mass is the
starting point of the theory, not its consequence (Doc.~365).

\subsection{The Zitterbewegung (1930, 1990)}

\textbf{What it contained.} Schrödinger found in the Dirac equation a rapid oscillation of the
electron's position with angular frequency $2mc^2/\hbar$ \QM~\cite{Schroedinger1930}. Hestenes
interpreted it as an internal circular motion and explicitly called the electron a clock
\QM~\cite{Hestenes1990}.

\textbf{What was missing.} The Zitterbewegung remained a peculiarity of single-particle Dirac
theory, regarded in the field picture as interference of positive- and negative-energy parts; no
geometric direction in which it takes place was assigned to it.

\textbf{Placement.} The corpus has not drawn on the Zitterbewegung so far. It belongs in this
series because it shows the same reciprocity of mass and rate in another language. An
identification with the circle direction $S^1_m$ is not claimed here; the factor~2 in the frequency
and the role of the negative-energy parts would first have to be clarified \SM.

\subsection{Pauli's objection to a time operator (1933)}

\textbf{What it contained.} Pauli showed that no self-adjoint time operator can exist for a
Hamiltonian that is bounded from below \QM~\cite{Pauli1933}. For decades this was taken as a reason
to keep time out of the state space.

\textbf{What was missing.} The objection presupposes open time. Compact time behaves like an angle
variable; its conjugate spectrum is discrete, as for angular momentum conjugate to an angle. This
gap in the argument was not exploited consistently.

\textbf{Where it fits in the corpus.} Doc.~307 and Doc.~334 record that Pauli's objection does not
apply to the rolled-up view: for compact time the Pauli condition is not satisfied, and the
discreteness of the modes becomes content rather than obstacle.

\subsection{Relational and timeless approaches (1983 onwards)}

\textbf{What it contained.} Page and Wootters obtain time from correlations in a static overall
state \QM~\cite{PageWootters1983}; the Wheeler-DeWitt equation and Rovelli's relational time
proceed similarly. Kaluza-Klein variants with compact time and Bars' two-time physics have been
studied \QM~\cite{Bars2001}.

\textbf{What was missing.} The mainstream takes time \emph{out} as something fundamental. The
opposite route, taking time \emph{in} as a compact dimension and obtaining the mass modes from it,
has hardly been pursued.

\textbf{Where it fits in the corpus.} Doc.~307 draws exactly this line: shared is the rejection of an
external parameter time; different is the positive proposal, compact geometric time instead of
relational or thermal time.

% ============================================================
\section{Mass as structure}
% ============================================================

\subsection{Wigner: mass as a Casimir invariant (1939)}

\textbf{What it contained.} Wigner classified particles as irreducible representations of the
Poincaré group; mass is a Casimir label there, hence already a structural eigenvalue and not an
attached number \EM~\cite{Wigner1939}.

\textbf{What was missing.} The classification does not say which values this label takes; the mass
values remain inputs.

\textbf{Where it fits in the corpus.} Doc.~307 cites Wigner as the reference showing that mass as a
structural quantity is well familiar; Doc.~339 uses the Wigner classification as a standard step in
separating the massive and massless sectors. In FFGFT the mass values themselves come from the
compact direction, not from the Poincaré group.

\subsection{The Casimir value of $SU(3)$}

\textbf{What it contained.} The quadratic Casimir operator of the fundamental representation of
$SU(N)$ has the value $(N^2-1)/(2N)$, for $SU(3)$ therefore $4/3$ \EM~\cite{Georgi1999}. The value is
textbook material of colour dynamics.

\textbf{What was missing.} Nobody had a reason to connect this value with a mode count of a compact
geometry.

\textbf{Where it fits in the corpus.} In Doc.~324, $\xi=C_2(SU(3)_{\mathrm{fund}})/N_{\mathrm{Fourier}}
=(4/3)/10^4=4/30000$; the formula is \BM{} there, the Fourier mode count \KM. The same value follows
in Doc.~338 from group orders of finite fields (section~\ref{sec:galois390}).

\subsection{Heaviside-Lorentz units (1893)}

\textbf{What it contained.} Heaviside introduced rationalised units in which the factors $4\pi$
disappear from the field equations \QM~\cite{Heaviside1893}. With $\hbar=c=\varepsilon_0=1$,
$\alpha=1$ becomes a choice of units.

\textbf{What was missing.} The choice of units remained a computational convenience; that $\hbar$
and $c$ are mere conversion factors became self-evident only with the practice of natural units and
the SI reform of 2019 (Doc.~312).

\textbf{Where it fits in the corpus.} Doc.~365 sets up the theory in this basic form: $\alpha=1$ is not
a setting but a choice of units, and the electromagnetic quantities become derived forms of mass
and time. The familiar number $137$ arises only in the transition to SI units; its agreement with
the measured value depends on the anchor $E_0$ (Doc.~338).

% ============================================================
\section{Geometry and algebra}
% ============================================================

\subsection{Orbifold compactification (1985)}

\textbf{What it contained.} Dixon, Harvey, Vafa and Witten compactified string theory on orbifolds,
tori divided by a finite group; $T^6/\mathbb{Z}_3$ was one of the first examples
\QM~\cite{DHVW1985}. The fixed points of the group carry sectors of their own there, and the number
of generations becomes a topological quantity.

\textbf{What was missing.} Six additional spatial dimensions and a vast number of possible
compactifications; time stays outside. Doc.~311 applies the corpus's two conditions to string
theory: the compact directions are exemplarily compact, all information is rebooked into conserved
quantum numbers, and what remains open is what the directions carry.

\textbf{Where it fits in the corpus.} The carrier is $T^4/\mathbb{Z}_3$ on the $D_4$ lattice, four
directions instead of six additional ones, with the time-mass direction among them (Doc.~330,
Doc.~314). The $\mathbb{Z}_3$ acts freely on all non-zero momenta; in position space the orbifold has
nine fixed points (R131 in Doc.~190). Torus, $\mathbb{Z}_3$ and $D_4$ are settings \SetzM{}
(Doc.~365).

\subsection{Finite fields and the Frobenius automorphism}
\label{sec:galois390}

\textbf{What it contained.} The finite fields $GF(p^n)$, their cyclic multiplicative groups and the
Frobenius automorphism $x\mapsto x^p$ have been fully described since Galois \EM~\cite{Lidl1997}.
They were not connected with elementary particle physics.

\textbf{What was missing.} A physical structure in which the prime~3 and its extensions occur
naturally.

\textbf{Where it fits in the corpus.} The $\mathbb{Z}_3$ of the orbifold supplies this structure. The
Frobenius automorphism of $GF(9)$ corresponds to the sector pairing $k\leftrightarrow-k$ \BM{}
(Doc.~336). Charge quantisation follows from the Legendre symbol modulo~13 \BM{} (Doc.~346), the
three generations arrange themselves in the Frobenius orbits of $GF(27)^*$ \BM{} (Doc.~348), and the
identity $(r_\mu/r_e)^2/\xi=|GF(9)^*|^2\cdot5^2\cdot|GF(27)|=43200$ connects the lepton ladder with
the group orders \KM{} (Doc.~338). The hypercharge of the charged leptons is still open here (R146
in Doc.~190).

\subsection{The Koide formula (1981--1983)}

\textbf{What it contained.} Koide found that the three charged lepton masses satisfy the relation
$Q=(m_e+m_\mu+m_\tau)/(\sqrt{m_e}+\sqrt{m_\mu}+\sqrt{m_\tau})^2=2/3$ to about $10^{-5}$
\QM~\cite{Koide1983}. For more than forty years the formula remained unexplained, a curiosity of
phenomenology.

\textbf{What was missing.} A reason for the value $2/3$ and for the phase that the formula requires
in its cosine form.

\textbf{Where it fits in the corpus.} The Koide phase is $\theta=2/9$ \BM, the same quotient as the
squared modulus of an icosahedral projection (Doc.~367). The amplitude $d/c=\sqrt2$ has the right
value but is not yet derived from the geometry and is treated as a setting \SM{} (Docs.~353, 367;
R150 in Doc.~190). A further puzzle piece fits in here afterwards: Foot showed in 1994 that $Q=2/3$
means an angle of $45^\circ$ between the vector of the $\sqrt{m_k}$ and the axis $(1,1,1)$
\QM~\cite{Foot1994}. In the language of the corpus this is the equipartition of the mass sum between
the $\mathbb{Z}_3$-symmetric mode and the pair of non-trivial modes \BM{} (Doc.~353). The open
question is thus no longer why $\sqrt2$ in particular, but why the mass structure divides half and
half between these two parts. The exact pair $(\sqrt2,\,2/9)$ is excluded for $m_\mu/m_e$
(Doc.~369). The formula is thus placed, but not completely closed.

% ============================================================
\section{Why the pieces did not come together}
% ============================================================

Three reasons run through all the sections.

\textbf{The pieces lay in different fields.} Kaluza-Klein belongs to gravitation and later to
string theory, the internal clock and the Zitterbewegung to quantum mechanics, Wigner and the
Casimir value to group theory, finite fields to number theory, the Koide formula to phenomenology.
Whoever works in one of these fields rarely has reason to read the others as parts of the same
picture.

\textbf{Each piece was classified as a tool or a curiosity.} The Compton time counted as a
conversion, Heaviside's units as a convenience, the Zitterbewegung as an artefact, the Koide formula
as a coincidence. The missing step was the same each time as with Einstein: to take a known relation
not as a derived quantity but as a principle (Doc.~312).

\textbf{Two switches were set wrongly.} Pauli's objection kept time out of the state space, and
Kaluza-Klein rolled up a spatial instead of the temporal direction. Both switches could have been
reset with the same insight: compact time escapes the objection, and it delivers the mass modes
without a free radius.

There is also a methodological point. That the pieces fit into the picture is a structural
agreement, not a numerical one. The more complex the quantum numbers, the closer one comes to the noise floor, where
a matching combination can always be found; no status is assigned there (Doc.~342, Doc.~343).
Scheme-dependent quantities of the Standard Model such as $\Lambda_{\mathrm{QCD}}$ are excluded from
any claim of derivation. The puzzle pieces therefore hold only where they concern simple,
convention-free ratios, and that is exactly where the places in this document lie.

% ============================================================
\section{Overview}
% ============================================================

\begin{center}
\small
\begin{tabular}{@{}>{\raggedright\arraybackslash}p{3.0cm}>{\raggedright\arraybackslash}p{1.7cm}>{\raggedright\arraybackslash}p{4.1cm}>{\raggedright\arraybackslash}p{3.8cm}@{}}
\toprule
Building block & Year & What was missing & FFGFT in the corpus \\
\midrule
Kaluza-Klein & 1921/26 & spatial direction, free radius, unobserved tower & Doc.~330, 270, 307, 311, 318 \\
De Broglie, Compton time & 1924/2013 & clock only as a derived scale & Doc.~312, 314, 365 \\
Zitterbewegung & 1930/1990 & no geometric direction & not used so far \SM \\
Pauli's objection & 1933 & presupposes open time & Doc.~307, 334 \\
Relational time, two-time physics & 1983 on & time taken out instead of in & Doc.~307 \\
Wigner & 1939 & mass values remain inputs & Doc.~307, 339 \\
Casimir value $SU(3)$ & textbook & no link to modes & Doc.~324, 338 \\
Heaviside-Lorentz & 1893 & only a computational convenience & Doc.~365, 312 \\
Orbifold & 1985 & six spatial directions, time outside & Doc.~330, 314, 311; R131 \\
Finite fields, Frobenius & Galois & no physical structure & Doc.~336, 338, 346, 348 \\
Koide, Foot & 1981--83, 1994 & no reason for $2/3$ and the phase & Doc.~367, 353, 369; R150 \\
\bottomrule
\end{tabular}
\end{center}

% ============================================================
\section{Conclusion}
% ============================================================

FFGFT connects to established physics and mathematics at eleven places, and at each of them it could
have been found. Kaluza-Klein had the compact direction, de Broglie the clock, Pauli the objection
whose gap is compact time, Wigner mass as structure, group theory the value $4/3$, string theory the
orbifold, number theory the finite fields, and Koide the puzzle they help to solve together. FFGFT is not
built from these pieces. It arose independently from time as a compact direction carrying the mass
modes, and from the algebra that follows from it with the one parameter $\xi$. The building blocks
were found in the literature only afterwards, and they fit in. That does not speak against the
theory but for it: pieces that arose independently of one another fit into a picture that was
drawn without them.

\vspace{1em}
\noindent\textit{Verification script:} \texttt{2/python/Dok390\_Skripte/pruef\_390\_puzzleteile.py}
--- 59/59.

% ============================================================
\section*{Tool and method}
\addcontentsline{toc}{section}{Tool and method}
% ============================================================

Claude (Anthropic) was used as a tool for: compiling the references in the corpus, writing and
running the verification script, and working out the \LaTeX{} sources according to the author's
specifications. The verification script backs every corpus reference in this document with a text
check in the named chapter file; if a check fails, the reference is not used.

\begin{thebibliography}{99}

\bibitem{Kaluza1921}
Th.~Kaluza, \textit{Zum Unitätsproblem der Physik},
Sitzungsber.\ Preuß.\ Akad.\ Wiss.\ Berlin (1921) 966--972.

\bibitem{Klein1926}
O.~Klein, \textit{Quantentheorie und fünfdimensionale Relativitätstheorie},
Z.\ Phys.\ \textbf{37} (1926) 895--906.

\bibitem{deBroglie1925}
L.~de~Broglie, \textit{Recherches sur la théorie des quanta},
Ann.\ Phys.\ (Paris) 10e sér.\ \textbf{3} (1925) 22--128 (thesis 1924).

\bibitem{Lan2013}
S.-Y.~Lan, P.-C.~Kuan, B.~Estey, D.~English, J.~M.~Brown, M.~A.~Hohensee, H.~Müller,
\textit{A clock directly linking time to a particle's mass},
Science \textbf{339} (2013) 554--557.

\bibitem{Schroedinger1930}
E.~Schrödinger, \textit{Über die kräftefreie Bewegung in der relativistischen Quantenmechanik},
Sitzungsber.\ Preuß.\ Akad.\ Wiss., Phys.-Math.\ Kl.\ \textbf{24} (1930) 418--428.

\bibitem{Hestenes1990}
D.~Hestenes, \textit{The zitterbewegung interpretation of quantum mechanics},
Found.\ Phys.\ \textbf{20} (1990) 1213--1232.

\bibitem{Pauli1933}
W.~Pauli, \textit{Die allgemeinen Prinzipien der Wellenmechanik},
in: Handbuch der Physik, vol.~24/1, Springer, Berlin 1933.

\bibitem{PageWootters1983}
D.~N.~Page, W.~K.~Wootters, \textit{Evolution without evolution: Dynamics described by stationary
observables}, Phys.\ Rev.\ D \textbf{27} (1983) 2885--2892.

\bibitem{Bars2001}
I.~Bars, \textit{Survey of two-time physics},
Class.\ Quantum Grav.\ \textbf{18} (2001) 3113--3130.

\bibitem{Wigner1939}
E.~Wigner, \textit{On unitary representations of the inhomogeneous Lorentz group},
Ann.\ Math.\ \textbf{40} (1939) 149--204.

\bibitem{Georgi1999}
H.~Georgi, \textit{Lie Algebras in Particle Physics}, 2nd ed., Westview Press, 1999.

\bibitem{Heaviside1893}
O.~Heaviside, \textit{Electromagnetic Theory}, vol.~1, The Electrician, London, 1893.

\bibitem{DHVW1985}
L.~Dixon, J.~A.~Harvey, C.~Vafa, E.~Witten, \textit{Strings on orbifolds},
Nucl.\ Phys.\ B \textbf{261} (1985) 678--686.

\bibitem{Lidl1997}
R.~Lidl, H.~Niederreiter, \textit{Finite Fields}, 2nd ed., Cambridge University Press, 1997.

\bibitem{Foot1994}
R.~Foot, \textit{A note on Koide's lepton mass relation}, arXiv:hep-ph/9402242 (1994).

\bibitem{Koide1983}
Y.~Koide, \textit{A fermion-boson composite model of quarks and leptons},
Phys.\ Lett.\ B \textbf{120} (1983) 161--165.

\end{thebibliography}
